\mxunit[4]{terminale-c04}{Limites de fonctions}\label{ch:3}

\mxsubtitle{S21, S23, S24, S26}

\begin{keypoint}[Objectifs]Calculer des limites en un point et à l'infini ; lever une indétermination ; interpréter une asymptote.\end{keypoint}

\section{Notes de cours}
\subsection{Sens d'une limite}
$f(x)\to\ell$ lorsque $x\to+\infty$ signifie que $f(x)$ peut être rendu aussi proche que souhaité de $\ell$ pour tous les $x$ suffisamment grands. Pour une limite en $a$, on considère des $x$ du domaine, proches de $a$ et distincts de $a$. Les limites à gauche et à droite se notent $\lim_{x\to a^-}f(x)$ et $\lim_{x\to a^+}f(x)$. Lorsque le domaine possède des points des deux côtés de $a$, une limite bilatérale existe si et seulement si les deux limites latérales existent et sont égales. À une extrémité du domaine, on étudie la limite du seul côté accessible.
\begin{keypoint}[Interprétation graphique]
Si $f(x)\to\ell\in\R$ lorsque $x\to+\infty$ ou $-\infty$, la droite $y=\ell$ est une asymptote horizontale dans cette direction. Si $f(x)\to\pm\infty$ lorsque $x\to a$ d'un côté, la droite $x=a$ est une asymptote verticale.
\end{keypoint}
Une courbe peut couper une asymptote horizontale. La limite ne donne pas nécessairement la position de la courbe par rapport à l'asymptote : étudier le signe de $f(x)-\ell$.
\subsection{Références et opérations}
En $+\infty$, $x^m\to+\infty$ pour $m\ge1$, $\frac{1}{x^m}\to0$, $\sqrt x\to+\infty$. En $-\infty$, le signe de $x^m$ dépend de la parité de $m$. Les règles de somme, produit, quotient et comparaison sont celles des suites, avec les mêmes formes indéterminées.
\paragraph{Exemple.} Pour $f(x)=\frac{2x+1}{x-3}$, écrire $f(x)=2+\frac{7}{x-3}$. Les limites en $\pm\infty$ valent $2$. En $3^-$, $f\to-\infty$ ; en $3^+$, $f\to+\infty$. Les asymptotes sont $x=3$ et $y=2$.
\paragraph{Conjugaison.} Pour $x>0$, $\sqrt{x^2+4x}-x=\frac{4}{\sqrt{1+\frac{4}{x}}+1}\to2$. Pour $x<0$, attention : $\sqrt{x^2}=|x|=-x$.
\subsection{Exponentielle et croissance comparée}
\[\lim_{x\to+\infty}e^x=+\infty,\quad\lim_{x\to-\infty}e^x=0,\quad
\lim_{x\to+\infty}\frac{e^x}{x^m}=+\infty\quad(m\in\N).\]
On en déduit $x^me^{-x}\to0$ en $+\infty$.
\paragraph{Démonstration des limites de l'exponentielle.} L'étude de $e^x-x-1$ donne $e^x\ge1+x$. Par comparaison, $e^x\to+\infty$ en $+\infty$. En posant $t=-x\to+\infty$, $e^x=\frac{1}{e^t}\to0$.
\paragraph{Démonstration d'une croissance comparée.} Pour tout entier $k\ge0$ et $x\ge0$, $e^x\ge \frac{x^k}{k!}$. Pour $k=0$, $e^x\ge1$. Si l'inégalité est acquise à l'ordre $k$, la dérivée de $e^x-\frac{x^{k+1}}{(k+1)!}$ est $e^x-\frac{x^k}{k!}\ge0$ ; sa valeur en $0$ est positive. L'hérédité suit. En choisissant $k=m+1$, $\frac{e^x}{x^m}\ge \frac{x}{(m+1)!}\to+\infty$.
\begin{method}[Choisir une transformation]
Quotient de polynômes : factoriser les plus grandes puissances. Différence de racines : quantité conjuguée. Quotient près d'un zéro : factoriser et étudier les signes. Exponentielle contre puissance : croissance comparée, après changement de variable si nécessaire.
\end{method}
\subsection{Comparaison et composition : conditions d'emploi}
Si $g(x)\le f(x)\le h(x)$ près du point ou de l'infini étudié, et si $g,h$ ont la même limite réelle $\ell$, alors $f(x)\to\ell$ (théorème des gendarmes). Si $g\le f$ et $g\to+\infty$, alors $f\to+\infty$. Pour une limite $-\infty$, on utilise une majoration par une fonction tendant vers $-\infty$.
Si $u(x)\to a$ et $v$ est continue en $a$, alors $v(u(x))\to v(a)$, lorsque la composée est définie. Par exemple $\sqrt{1+\frac{1}{x}}\to1$ en $+\infty$. Pour une limite infinie de la fonction extérieure, vérifier de quel côté $u(x)$ approche la borne : $\frac{1}{u(x)}$ n'a pas le même comportement selon le signe de $u(x)$.


\mxnewpage
\section{Activité de recherche}\textbf{Une courbe proche d'une droite --- 30 min.}\par On considère la fonction $f$ définie sur $\R\setminus\{1\}$ par $f(x)=\frac{3x-2}{x-1}$. On note $\mathcal C_f$ sa courbe dans un repère du plan.
\begin{enumerate}\item Montrer que, pour tout $x\ne1$, $f(x)=3+\frac1{x-1}$.
\item Déterminer les limites de $f$ en $-\infty$, en $1^-$, en $1^+$ et en $+\infty$.
\item Étudier le signe de $f(x)-3$, puis préciser la position de $\mathcal C_f$ par rapport à la droite d'équation $y=3$.
\item Déterminer un réel $A>1$ tel que, pour tout $x>A$, $|f(x)-3|<\num{0,01}$.
\item La courbe $\mathcal C_f$ coupe-t-elle son asymptote horizontale ? Justifier et expliquer pourquoi une fenêtre graphique trop large peut donner une impression trompeuse.\end{enumerate}

\mxnewpage
\section{Exercices progressifs}

\mxexerciseitem{1}{Entraînement}Déterminer $\displaystyle\lim_{x\to+\infty}\frac{5x^3-x}{2x^3+7}$ en justifiant.

\mxexerciseitem{2}{Entraînement}Soit $f$ la fonction définie sur $\R\setminus\{2\}$ par $f(x)=\frac1{(x-2)^2}$. Déterminer ses limites en $2^-$, en $2^+$, en $-\infty$ et en $+\infty$.

\mxexerciseitem{3}{Entraînement}Déterminer $\displaystyle\lim_{x\to1}\frac{x^2-1}{x-1}$ en précisant pourquoi la simplification utilisée est valable au voisinage de $1$.

\mxexerciseitem{4}{Entraînement}Déterminer $\displaystyle\lim_{x\to+\infty}x^3e^{-2x}$. Citer le résultat de croissance comparée utilisé.

\mxexerciseitem{5}{Entraînement}Déterminer $\displaystyle\lim_{x\to-\infty}\left(\sqrt{x^2+1}+x\right)$. Justifier notamment le traitement de $\sqrt{x^2}$.

\section{Travail dans le manuel ou une annale}\textbf{Barbazo, p. 96, exercices 26 et 29.} En S23 : exercice 29, pour relier limites et asymptotes. À la maison après S24 : exercice 26, pour construire trois contre-exemples à une règle de calcul abusive sur les limites.

\mxnewpage
\section{Devoir à la maison}\textbf{DM3 --- Une famille de courbes --- 45 min.}\par Pour chaque réel $a$, on considère la fonction $f_a$ définie par $f_a(x)=\frac{ax^2+1}{x^2+1}$. On note $\mathcal C_a$ sa courbe représentative.
\begin{enumerate}\item Déterminer l'ensemble de définition de $f_a$, puis ses limites en $-\infty$ et en $+\infty$.
\item Montrer que $f_a(x)-a=\frac{1-a}{x^2+1}$. Selon la valeur de $a$, préciser la position de $\mathcal C_a$ par rapport à la droite d'équation $y=a$.
\item Étudier les variations de $f_a$ en distinguant les cas $a<1$, $a=1$ et $a>1$. Préciser les éventuels extrema et où ils sont atteints.
\item On prend $a=3$. Résoudre, pour $x\ge0$, l'inéquation $|f_3(x)-3|\le\num{0,02}$.\end{enumerate}

\ifmxteacher\mxnewpage\section{Conduite et corrigés}\begin{teachernote}[Repères pour le professeur]En S21, relier chaque limite à une lecture graphique. En S23, exiger les signes du dénominateur à gauche et à droite. En S24, faire choisir la transformation avant de calculer : factorisation, conjugaison, développement local non utilisé. En S26, séparer la démonstration des croissances comparées de leur utilisation. Les asymptotes obliques ne sont pas requises ici.\end{teachernote}

\subsection{Corrigé de l'activité}\begin{enumerate}
\item Comme $3x-2=3(x-1)+1$, pour $x\ne1$, $f(x)=3+\frac{1}{x-1}$.
\item En $\pm\infty$, $\frac{1}{x-1}\to0$ : les limites valent $3$. En $1^-$, $x-1\to0^-$, donc $f(x)\to-\infty$ ; en $1^+$, $x-1\to0^+$, donc $f(x)\to+\infty$.
\item Le signe de $f(x)-3=\frac{1}{x-1}$ est négatif pour $x<1$, positif pour $x>1$. La courbe est donc respectivement au-dessous et au-dessus de l'asymptote $y=3$.
\item Pour $x>1$, $|f(x)-3|<\num{0,01}$ équivaut à $\frac{1}{x-1}<\num{0,01}$, puis $x>101$. On peut choisir $A=101$.
\item Non : $f(x)-3=\frac{1}{x-1}$ n'est jamais nul. La courbe peut paraître confondue avec l'asymptote à la précision de l'écran, mais elle ne la rencontre pas.
\end{enumerate}

\paragraph{E1.}Pour $x\ne0$, le quotient vaut $\frac{5-\frac{1}{x^2}}{2+\frac{7}{x^3}}$. En $+\infty$, les puissances inverses tendent vers $0$ et le dénominateur vers $2\ne0$. La limite est $\frac{5}{2}$.

\paragraph{E2.}En $2$, $(x-2)^2\to0$ par valeurs strictement positives des deux côtés ; son inverse tend donc vers $+\infty$. En $\pm\infty$, le carré tend vers $+\infty$ et son inverse vers $0$.

\paragraph{E3.}Pour $x\ne1$, $x^2-1=(x-1)(x+1)$, donc le quotient est égal à $x+1$. Sa limite en $1$ est donc $2$. Cette simplification calcule la limite ; elle ne rend pas la fonction initiale définie en $1$.

\paragraph{E4.}Posons $t=2x\to+\infty$. Alors $x^3e^{-2x}=\frac{t^3}{8e^t}\to0$, par croissance comparée de l'exponentielle et de la puissance.

\paragraph{E5.}Le conjugué donne $\sqrt{x^2+1}+x=\frac{1}{\sqrt{x^2+1}-x}$. Lorsque $x\to-\infty$, les deux termes $\sqrt{x^2+1}$ et $-x$ tendent vers $+\infty$ ; le dénominateur tend donc vers $+\infty$. La limite vaut $0$, par valeurs positives. Pour $x<0$, $\sqrt{x^2}=|x|=-x$ : la somme $\sqrt{x^2+1}+x$ est donc une forme $\infty-\infty$, d'où le recours au conjugué.

\subsection{Corrigé du DM3}\begin{enumerate}
\item Comme $x^2+1>0$ sur $\R$, le domaine est $\R$. En divisant par $x^2$ pour $x\ne0$, $f_a(x)=\frac{a+\frac{1}{x^2}}{1+\frac{1}{x^2}}\to a$ en $\pm\infty$.
\item $f_a(x)-a=\frac{ax^2+1-a(x^2+1)}{x^2+1}=\frac{1-a}{x^2+1}$. Pour $a<1$, la courbe est au-dessus de $y=a$ ; pour $a>1$, elle est en dessous ; pour $a=1$, elle est confondue avec cette droite.
\item Par quotient, $f'_a(x)=\frac{[2ax(x^2+1)-2x(ax^2+1)]}{(x^2+1)^2}=\frac{2x(a-1)}{(x^2+1)^2}$.
Si $a>1$, le signe est celui de $x$ : décroissance sur $]-\infty,0]$, croissance sur $[0,+\infty[$ ; minimum $f_a(0)=1$ et pas de maximum (la borne $a$ n'est pas atteinte). Si $a<1$, les variations sont inversées : maximum $1$ et pas de minimum. Pour $a=1$, la fonction est constante, de valeur $1$.
\item Pour $a=3$ et $x\ge0$, l'inégalité devient $\frac{2}{x^2+1}\le\num{0,02}$. Le dénominateur est positif, donc elle équivaut à $x^2+1\ge100$, soit $x\ge\sqrt{99}$.
\end{enumerate}

\subsection{Corrigé du travail dans le manuel ou une annale}\paragraph{Exercice 26.} Lorsque $x\to+\infty$, les fonctions $f(x)=x$ et $g(x)=x$ tendent toutes deux vers $+\infty$, mais leur différence est nulle. Avec $f(x)=2x$, $g(x)=x$, la différence tend vers $+\infty$ ; en échangeant $f$ et $g$, elle tend vers $-\infty$. Ces exemples établissent le caractère indéterminé de la différence.
\paragraph{Exercice 29.} Le dénominateur s'annule en $\frac32$, donc $D_f=\R\setminus\{\frac32\}$. On écrit
\[f(x)=\frac{5-2x}{2x-3}=-1+\frac2{2x-3}.\]
Ainsi $f(x)\to-1$ aux deux infinis. En $\frac32$, le numérateur du second terme est positif et son dénominateur tend vers $0$ : la limite à gauche est $-\infty$, celle à droite est $+\infty$. Les asymptotes ont pour équations $y=-1$ et $x=\frac32$.

\fi

